\documentclass[12pt]{article}
\usepackage[utf8]{inputenc}
\usepackage[T1]{fontenc}
\usepackage{hyperref}
\title{Karawa Language Resources for K-12 Curriculum Development}
\author{Compiled by Claude Sonnet 5.5 (claude-sonnet-5-5) for Andr\'as Kornai}
\date{}
\begin{document}
\maketitle

\section{Introduction}
Karawa (also Bulawa; ISO 639-3 xrw) is an endangered Sepik language of the Ram group, together with Awtuw and Pouye. It is spoken in the village of Pulwa (Bulawa) in East Wapei Rural LLG, Sandaun Province, Papua New Guinea. Reported speaker numbers are 63 (Wikipedia, 2003) and 44 (Verbix). Speakers are shifting to the closely related Pouye. The search found little on grammar, documentation or recent fieldwork.

\section{Primary Educational Resources}
No primers, textbooks or school programmes for Karawa were found.

\section{Linguistic and Reference Materials}
Wikipedia article: \url{https://en.wikipedia.org/wiki/Karawa_language}. Related group: \url{https://en.wikipedia.org/wiki/Ram_languages}. Glottolog: \url{https://glottolog.org/resource/languoid/id/kara1495.bigmap.html}. ASJP wordlist: \url{https://asjp.clld.org/languages/KARAWA}. Verbix endangered-language page: \url{https://docs.verbix.com/EndangeredLanguages/Karawa}.

\section{Religious and Community Resources}
Joshua Project listing: \url{https://joshuaproject.net/languages/xrw}.

\section{Digital and Audio Resources}
Global Recordings Network has a language page for xrw: \url{https://globalrecordings.net/en/language/xrw}. I did not verify what recordings it holds.

\section{Curriculum Development Context}
With so few speakers and shift toward Pouye, any curriculum would be a community-led documentation and heritage effort. It would start from a basic wordlist and recordings, and would probably be coordinated with Pouye or Tok Pisin programmes in the area.
\end{document}